(i)$\Rightarrow$(ii): For any ideal $I$, its radical is the intersection of the primes containing it. Replacing each prime by its intersection of maximal ideals gives $\sqrt I=\bigcap_{\mathfrak{m}\supseteq I,\ \mathfrak{m}\text{ maximal}}\mathfrak{m}$. Passing to $A/I$ proves (ii).

(ii)$\Rightarrow$(i): Apply (ii) to $A/\mathfrak{p}$. Its nilradical is zero, so $\mathfrak{p}$ is the intersection of the maximal ideals containing it.

(i)$\Rightarrow$(iii): If $\mathfrak{p}$ is not maximal, all maximal ideals containing it contain it strictly. Thus the intersection of all primes strictly containing it lies in their intersection, which is $\mathfrak{p}$.

(iii)$\Rightarrow$(i): Otherwise pass to $A/\mathfrak{p}$ for a prime not expressible as an intersection of maximal ideals. Condition (iii) passes to quotients, and we now have a domain with $J(A)\neq0$. Choose $0\neq f\in J(A)$ and contract a maximal ideal of $A_f$ to $\mathfrak{q}\subseteq A$. Then $f\notin\mathfrak{q}$, and every prime strictly containing $\mathfrak{q}$ contains $f$. Since all maximal ideals of $A$ contain $f$, the prime $\mathfrak{q}$ is not maximal. The intersection of the primes strictly containing it therefore contains $f\notin\mathfrak{q}$, contradicting (iii).
